\documentclass[11pt]{article}

\usepackage[margin=1in]{geometry}
\usepackage{amsmath}
\usepackage{booktabs}
\usepackage{hyperref}
\usepackage{xcolor}

\title{Methodological Summary: VIS Bilevel Evaluator-Prompt Evolution}
\author{}
\date{}

\begin{document}
\maketitle

\section{Purpose}

This experiment studies a bilevel optimization design for insight discovery on the
IEEE VIS benchmark. The inner loop evolves candidate insight-generation programs.
The outer loop evolves the evaluator system prompt when the inner loop appears
saturated. The main design change is a \emph{latest-only} objective: after an
evaluator prompt has evolved, new generator candidates are optimized under the
latest active evaluator prompt \(p_t\), while the original evaluator prompt
\(p_0\) is retained only as diagnostic context.

The goal is to reduce saturation caused by optimizing against a fixed rubric.
Instead of forcing all future search behavior to remain close to \(p_0\), the
outer loop searches for evaluator prompts that are textually diverse and still
locally stable relative to the current active evaluator. This creates a moving
local trust region:
\[
  p_t \rightarrow p_{t+1},
\]
rather than a permanent global anchor:
\[
  p_0 \rightarrow p_t.
\]

\section{Repository Layout}

The experiment assets are consolidated under:
\[
\texttt{experiments/vis50\_bilevel\_latest\_only/}.
\]

\begin{table}[h]
\centering
\begin{tabular}{ll}
\toprule
Path & Purpose \\
\midrule
\texttt{run.sh} & Reproduce the VIS 50-iteration run. \\
\texttt{analyze.py} & Generate interval figures and prompt-progression reports. \\
\texttt{methodological\_summary.tex} & This methodology document. \\
\texttt{results/} & Consolidated experiment outputs. \\
\bottomrule
\end{tabular}
\end{table}

The implementation lives primarily in:
\begin{itemize}
  \item \texttt{skydiscover/search/adaevolve/controller.py}
  \item \texttt{skydiscover/search/adaevolve/evaluator\_prompt\_evolution.py}
  \item \texttt{skydiscover/context\_builder/adaevolve/builder.py}
  \item \texttt{benchmarks/insight\_gen/VIS/evaluator.py}
\end{itemize}

\section{Inner Loop: Generator Optimization}

At each iteration, the AdaEvolve controller selects a parent program, builds a
prompt for the generator LLM, asks the LLM to produce a child program, executes
the child, and evaluates the resulting insight. The generator prompt is assembled
by the AdaEvolve context builder. It includes:

\begin{itemize}
  \item the parent program and selected context programs;
  \item evaluator feedback and judge evidence from the parent and nearby programs;
  \item the latest evaluator score \(J(c; p_t)\);
  \item the old evaluator score \(J(c; p_0)\), when available, as diagnostic
        context only;
  \item paradigm-breakthrough guidance when the search controller detects global
        stagnation.
\end{itemize}

The latest evidence block is produced by
\texttt{\_format\_latest\_evaluator\_feedback}. Its prompt explicitly instructs
the generator to use feedback paired with \(\texttt{latest\_combined\_score}\)
from \(p_t\) when deciding the next move. In latest-only mode, the score used by
the optimizer is:
\[
  S_{\mathrm{gen}}(c) = J(c; p_t).
\]
The original score
\[
  J(c; p_0)
\]
is still logged as \(\texttt{old\_combined\_score}\), but it is not part of the
generator selection objective.

\section{Score Saturation Monitor}

Evaluator-prompt evolution is not run at every generator iteration. It is
triggered by a best-so-far saturation monitor implemented as
\texttt{ScoreSaturationMonitor}. When a child candidate receives a valid score,
the monitor updates
\[
  B_i = \max(B_{i-1}, S_{\mathrm{gen}}(c_i)),
\]
where \(S_{\mathrm{gen}}\) is the active generator selection score. In
latest-only mode, this is \(J(c_i; p_t)\). The monitor stores the sequence of
\((\text{iteration}, B_i)\) pairs rather than the raw per-iteration scores.

Given a window size \(w\), the monitor examines the most recent \(w\)
best-so-far values and computes
\[
  \Delta_{\mathrm{window}}
    = B_{\mathrm{last}} - B_{\mathrm{first}}.
\]
It declares saturation when
\[
  \Delta_{\mathrm{window}} < \epsilon,
\]
where \(\epsilon\) is the configured saturation threshold. The VIS experiment
uses
\[
  w = 5,\qquad \epsilon = 0.005.
\]
Thus, an evaluator-prompt update is triggered only when the best score has
improved by less than \(0.005\) across the last five recorded candidate scores.
Because the monitor tracks best-so-far scores, a failed or zero-scoring
iteration does not by itself lower the monitored value; it contributes to a
trigger only by extending a period with no new best score.

The monitor also enforces a minimum spacing between triggers. In this
experiment,
\[
  \text{\texttt{evaluator\_prompt\_min\_interval}} = 5,
\]
matching the saturation window. After a trigger is marked, another trigger is
suppressed until at least five generator iterations have passed. This prevents
the outer evaluator loop from repeatedly firing on the same stagnant window.

The outer evaluator search uses a separate internal saturation monitor over
outer-loop best fitness. Its window is
\[
  \max(3,\lfloor \text{\texttt{outer\_max\_iterations}}/3 \rfloor),
\]
with threshold \(10^{-4}\). This can stop an individual outer prompt-search run
early when the best outer fitness stops improving, but it is separate from the
inner monitor that decides whether evaluator evolution should begin.

\subsection{Recent-Window Feedback Passed to the Outer Guide LLM}

The saturation monitor itself stores only the iteration number and the
best-so-far scalar score. It does not store the program identifier, insight text,
judge evidence, or evaluator feedback attached to each best-so-far point. When
the monitor triggers evaluator-prompt evolution, the controller constructs a
separate recent-window feedback block from the program database.

This feedback block is produced by \texttt{aggregate\_feedback}. It selects
programs satisfying three conditions:
\begin{itemize}
  \item the program has an integer \(\texttt{iteration\_found}\);
  \item the program was found no later than the current saturation iteration;
  \item the program has evaluator metrics.
\end{itemize}
The selected programs are sorted by descending \(\texttt{iteration\_found}\),
and only the most recent \(w\) programs are retained, where \(w\) is the same
window size used by the saturation monitor. In the VIS experiment, this means
the feedback block contains at most the five most recent scored programs.

For each retained program, the block includes the iteration, scalar scores
\texttt{combined\_score}, \texttt{old\_combined\_score}, and
\texttt{latest\_combined\_score}, plus available evaluator evidence fields:
\[
\begin{gathered}
  \text{\texttt{judge\_evidence}},\quad
  \text{\texttt{judge\_conclusion}},\quad
  \text{\texttt{old\_judge\_evidence}},\quad
  \text{\texttt{old\_judge\_conclusion}},\\
  \text{\texttt{latest\_judge\_evidence}},\quad
  \text{\texttt{latest\_judge\_conclusion}},
\end{gathered}
\]
as well as any evaluator error fields. The resulting text is truncated to
\texttt{evaluator\_prompt\_max\_feedback\_chars} before being passed to the
outer guide LLM.

The call path is:
\[
\text{\texttt{\_maybe\_evolve\_evaluator\_prompt}}
\rightarrow
\text{\texttt{aggregate\_feedback}}
\rightarrow
\text{\texttt{\_run\_bilevel\_outer}}
\rightarrow
\text{\texttt{run}}/\text{\texttt{run\_single\_shot}}
\rightarrow
\text{\texttt{\_apply\_operator}}.
\]
Thus the outer guide LLM sees recent evidence from the current plateau window
when proposing a new evaluator prompt. This is not a replay of the entire
best-score history: if the last new best occurred before the recent window, its
feedback is not necessarily included in this block.

\section{Outer Loop: Evaluator-Prompt Evolution}

The outer loop is triggered when recent inner-loop scores satisfy the saturation
monitor. 

\subsection{Probe and Canary Sets}

The outer loop uses two program sets:

\begin{itemize}
  \item \textbf{Probe set}: a stratified top/mid/low sample from the current
        database. This set scores outer-loop candidate prompts.
  \item \textbf{Canary set}: a persistent stratified top/mid/low sample frozen
        at the first outer-loop trigger. This set measures cumulative score drift.
\end{itemize}

The canary membership stays fixed, but its reference scores are refreshed against
the configured reference evaluator before each outer step. In latest-only mode,
the reference evaluator is:
\[
  p_{\mathrm{ref}} =
  \begin{cases}
    p_0, & \text{if no evolved prompt exists}, \\
    p_t, & \text{otherwise}.
  \end{cases}
\]

\subsection{Outer Fitness}


\section{Drift Control}

There are two drift controls:

\begin{enumerate}
  \item \textbf{Probe drift}: included in the outer fitness through
        \(\Delta_{\mathrm{local}}\).
  \item \textbf{Canary cumulative drift}: computed on the persistent canary set
        before installing an outer-loop winner.
\end{enumerate}

The configured mode for the VIS experiment is \(\texttt{soft}\). In soft mode,
drift is penalized and logged, but it does not reject a candidate by itself. If
\(\texttt{outer\_evaluator\_drift\_control\_mode}\) is set to \(\texttt{strict}\),
the same machinery becomes a hard gate using
\(\texttt{outer\_evaluator\_drift\_hard\_cap}\) and
\(\texttt{outer\_evaluator\_cumulative\_drift\_cap}\).

\section{Guide LLM Prompts}

There are two guide-LLM uses.

\subsection{Generator Guidance}

The inner generator sees the latest evaluator evidence through the search
guidance section emitted by the context builder. The key instruction is that
\(\texttt{latest\_combined\_score}\) from \(p_t\) drives the next move, while
\(p_0\) evidence is diagnostic.

\subsection{Evaluator-Prompt Guidance}

The outer loop uses guide-LLM operators:

\begin{itemize}
  \item \textbf{branch}: request a structurally different evaluator prompt;
  \item \textbf{refine}: request a small targeted prompt edit;
  \item \textbf{merge}: combine two prompt structures.
\end{itemize}

All operators preserve the evaluator JSON contract, the score dimensions, and
the integer 0--100 score range. They are asked to increase prompt diversity while
keeping the chart-grounded evaluator task stable.

\section{Experiment Configuration}

The reproducibility script defaults to:

\begin{table}[h]
\centering
\begin{tabular}{ll}
\toprule
Setting & Value \\
\midrule
Model for generation & \texttt{gpt-4o-mini} \\
Model for evaluation & \texttt{gpt-4o-mini} \\
Model for guidance & \texttt{gpt-4o-mini} \\
Iterations & 50 \\
Saturation window & 5 \\
Generator score mode & \texttt{latest\_only} \\
Outer evaluator mode & \texttt{adaevolve} \\
Outer baseline mode & \texttt{latest\_only} \\
Outer drift control & \texttt{soft} \\
Outer max iterations per trigger & 8 \\
Outer population size & 4 \\
Outer samples per eval & 2 \\
\bottomrule
\end{tabular}
\end{table}

\section{Telemetry and Analysis}

The run records inner-loop iteration stats in an
\(\texttt{adaevolve\_iteration\_stats\_*.jsonl}\) file. It records outer-loop
prompt-search telemetry in \(\texttt{outer\_evaluator\_stats.jsonl}\), including:

\begin{itemize}
  \item saturation iteration;
  \item outer iteration;
  \item operator;
  \item prompt diversity;
  \item discrimination;
  \item within variance;
  \item drift;
  \item drift-control mode;
  \item baseline mode;
  \item baseline prompt version;
  \item install or rejection outcome.
\end{itemize}

The analysis script generates:

\begin{itemize}
  \item \(\texttt{evaluator\_intervals.png}\): score intervals between evaluator
        prompt versions, with best-so-far under \(p_0\) as a dashed grey line;
  \item \(\texttt{evaluator\_interval\_summary.md}\): interval-level summary;
  \item \(\texttt{evaluator\_prompt\_progression\_local\_changes.md}\): local
        diffs for prompt changes.
\end{itemize}

\section{Interpretation}

This design intentionally relaxes the old global \(p_0\) constraint. If the
outer loop always optimizes for diversity while staying close to \(p_0\), prompt
evolution can become self-contradictory: the search asks for novelty but also
keeps pulling every candidate back to the original rubric. The latest-only design
instead asks each evaluator update to make a controlled local move from the
current prompt. This makes it easier for the evaluator to leave a saturated
region while still retaining measurable score-behavior continuity.

The tradeoff is that long-run alignment to \(p_0\) is no longer enforced by the
optimization objective. That is intentional for this experiment. The remaining
guardrail is local drift, plus full telemetry for \(p_0\) diagnostics so that
post-hoc analysis can show whether the evolved evaluator materially diverged from
the original rubric.

\end{document}
